\chapter{Geometry}
\label{ch:geometry}

\section{Triangle properties}
\label{sec:triprops}

For a triangle with corners $\vect{x}_0,\vect{x}_1,\vect{x}_2$ (in winding order) let
$\vect{a}=\vect{d}(\vect{x}_0,\vect{x}_1)$ and $\vect{b}=\vect{d}(\vect{x}_0,\vect{x}_2)$ be the
minimum-image edge vectors. Then
\begin{gather}
  A = \tfrac12\lvert\vect{a}\times\vect{b}\rvert,\qquad
  \vect{n} = \frac{\vect{a}\times\vect{b}}{\max(\lvert\vect{a}\times\vect{b}\rvert,10^{-15})},\\
  \vect{c} = \operatorname{wrap}\!\left(\vect{x}_0 + \tfrac13(\vect{a}+\vect{b})\right).
\end{gather}
\fn{MFTLGetTriangleProperties} computes them from the current positions;
\fn{MFTLComputeTriangleProperties} stores them in the owned triangles of an object
(\texttt{area}, \texttt{normal}, \texttt{center}). The cached values are what the ownership rule,
the integrals, the curvature kernels and the intersection tests read. Ghost triangles receive
their cached values from the owner at the next redistribution.

\section{Single-body integrals}

The integrals are collective (a sum over owned triangles, reduced over all ranks) and defined
for a body that encloses a volume by itself: closed, and compact on every periodic axis. On a
periodic domain every triangle is evaluated in continuous coordinates in the image of one
collectively agreed anchor vertex, so a body of any extent across seams is integrated exactly. A
wrapping or open body makes the volume and the moment of inertia abort collectively (the reason is
printed on rank 0) --- use a phase-boundary set instead.

\paragraph{Volume} (\fn{MFTLObjectVolume}):
\begin{equation}
  V = \frac16\sum_{T\ \text{owned}} \vect{x}_1\cdot(\vect{x}_2\times\vect{x}_3).
\end{equation}
Its sign follows the winding: $V>0$ for the outward convention.

\paragraph{Surface centre} (\fn{MFTLSurfaceCenter}): the area-weighted centroid of the surface
(\emph{not} the centroid of the enclosed volume),
\begin{equation}
  \vect{c}_S = \frac{\sum_T A_T\,(\vect{x}_1+\vect{x}_2+\vect{x}_3)}{3\sum_T A_T}.
\end{equation}
The function returns the total area. It is defined for open bodies too. On a periodic domain
the centre is a continuous coordinate and may lie outside the box.

\paragraph{Moment of inertia} (\fn{MFTLMomentOfInertia}): the \mftlone{} lumped-tetrahedron
approximation about a centre $\vect{c}$ with uniform density $\rho$. Every owned triangle and
$\vect{c}$ span a tetrahedron of volume $|V_T|$, whose mass is lumped at its centroid; with
$\vect{r}_T=\tfrac14(\vect{x}_1+\vect{x}_2+\vect{x}_3+\vect{c})-\vect{c}$,
\begin{equation}
  I_{ij} = \rho\sum_{T}|V_T|\left(\lvert\vect{r}_T\rvert^2\delta_{ij}-r_{T,i}\,r_{T,j}\right).
\end{equation}
This is not the exact tetrahedron covariance formula; it is kept for \mftlone{} parity. The
centre must be given in the same periodic image as the mesh, for example the surface centre.

\section{Local-frame kernels}

\texttt{mftl\_geometry.h} provides the kernels from which periodic-safe geometry is built:
\fn{MFTLVertexDelta} (minimum-image displacement between two vertices), \fn{MFTLEdgeVector} (of a
half-edge), \fn{MFTLTriangleUnwrapped} (the corners of a triangle in the image of its first
corner), \fn{MFTLVertexAngleSum} and \fn{MFTLVertexTurningAngle} (angle sums and the geodesic
turning angle $\pi-\sum\alpha$ of a boundary vertex), \fn{MFTLBoundaryNeighbours} (the loop
neighbours of a boundary vertex), \fn{MFTLBoundaryEdgeGeometry} (Figure~\ref{fig:conormal}) and
\fn{MFTLBoundaryEdgeCentroidDistance} (the in-surface distance $\delta$ from a boundary
triangle's centroid to the edge midpoint and the height $h_T=2A_T/L_e$; $\delta\ge h_T/3$).
\fn{MFTLFindThinRegionPairs} and \fn{MFTLFindProximityPairs} find stored triangle pairs that come
close (thin regions of one body; two bodies).

\section{Transforms}

\begin{lstlisting}[numbers=none]
void MFTLTranslateObject(MFTL_Object *obj, const double d[3]);
int  MFTLScaleObject(MFTL_Object *obj, const double s[3]);
int  MFTLRotateObject(MFTL_Object *obj, const double R[3][3], const double centre[3]);
void MFTLWrapObjectIntoDomain(MFTL_Object *obj);
\end{lstlisting}
The transforms move the \emph{owned} vertices and recompute the triangle properties; the ghosts
are refreshed afterwards with \fn{MFTLUpdateGhostPositions}. \fn{MFTLTranslateObject} is
rank-local and does not wrap: positions may leave the box with unchanged counters, and
\fn{MFTLWrapObjectIntoDomain} wraps them and counts the wraps. On a periodic domain scaling (per
axis, about the origin) and rotation (about \texttt{centre}) are collective: they act on
continuous coordinates relative to the body reference point, about the image of the origin or
centre nearest to it, and refuse (returning \texttt{MFTL\_ERROR}, body unchanged) a map that
would take a winding vector off the periodic lattice.

\section{Normals and curvature}

All vertex quantities are one-ring sums over the stored fan of an owned vertex. Their
preconditions: the fan is stored (G2, guaranteed after a redistribution), the ghost positions
are current, and the cached triangle properties are current.

\paragraph{Vertex normal} (\fn{MFTLComputeVertexNormals}): the incident triangle normals
weighted by $w=\sin\alpha/(\ell_0\ell_1)$, where $\alpha$ is the corner angle at the vertex and
$\ell_0,\ell_1$ are the lengths of the two edges at that corner, normalised to unit length.

\paragraph{Curvature} (\fn{MFTLComputeVertexCurvatures}): the cotangent Laplace--Beltrami
operator with the mixed Voronoi area $A_{\text{mix}}$ (Meyer et al.):
\begin{align}
  \vect{K}(v) &= \frac{1}{2A_{\text{mix}}}\sum_{j}\left(\cot\alpha_{j}+\cot\beta_{j}\right)(\vect{x}_v-\vect{x}_j),
  &\texttt{meancurvature} &= \lvert\vect{K}(v)\rvert \;(=2H),\\
  \texttt{gaussiancurvature} &= \frac{2\,(2\pi-\sum_{T}\theta_T)}{A_{\text{mix}}},
\end{align}
where $\alpha_j,\beta_j$ are the angles opposite the edge $vj$ and $\theta_T$ the corner angles at
$v$. On a sphere of radius $R$ the mean curvature is $2/R$ and the Gaussian curvature $2/R^2$
(the \mftlone{} convention; the remesher's curvature-adaptive edge length depends on it). The
mixed area of an obtuse triangle contributes $A/2$ at the obtuse corner and $A/4$ otherwise.
Single-vertex versions \fn{MFTLComputeVertexNormalAt} and \fn{MFTLComputeVertexCurvatureAt} work on
any stored vertex (owned or ghost) with a closed stored fan.

\paragraph{Boundaries.} A boundary vertex gets the angle-weighted normal of its open fan; its
Laplace--Beltrami curvature is not defined and it is flagged
\texttt{MFTL\_VFLAG\_CURV\_UNDEFINED} (never silently skipped).
\fn{MFTLComputeExtrapolatedCurvatures} fits a least-squares paraboloid over the open two-ring of
each owned boundary vertex (at least six points) and stores $\lvert2H\rvert$ in
\texttt{curv\_extrap} with \texttt{MFTL\_VFLAG\_CURV\_EXTRAPOLATED}.

\paragraph{Roughness} (\fn{MFTLComputeRoughness}, collective): returns the global average of the
geometric Laplacian $\lvert\vect{x}_v-\sum_j(\vect{x}_j/d_j)/\sum_j(1/d_j)\rvert$ and stores per
vertex the mean and sample variance of the dihedral roughness $1-\vect{n}_{T_0}\cdot\vect{n}_{T_1}$
of the incident edges (\texttt{roughness\_mean}, \texttt{roughness\_variance}). The remesher's
roughness-triggered smoothing uses them.

Values on ghost vertices are refreshed with \fn{MFTLUpdateGhostPositions}, which carries normal,
curvatures and area. The example of \S\ref{sec:ex-curvature} computes these quantities on a
distributed sphere.
